\documentclass[11pt]{article}
\usepackage[a4paper,margin=25mm]{geometry}
\usepackage{amsmath,amssymb,amsthm,hyperref}
\newtheorem{proposition}{Proposition}
\title{The mixed-integer source does not force its ninth input to be nonnegative}
\author{\(\mathbb Z\mathbb Q\) Diophantine Project}
\date{23 November 2026}
\begin{document}
\maketitle

Use the notation of the exact BDHMS coefficient recipe in
\texttt{finite\_wrapper\_transfer.tex}: the source is
\(Q(a,f,g,h,k,l,w,x,y,n)\), with eight integer inputs and one
nonnegative integer input.  The coding quantities are denoted by
\(B_0,N,S_0,T_0,R_0,X,Y\), to distinguish them from the final
relation-combining arguments \(S,T,R\).
The primary source is J.~Bayer, M.~David, M.~Ha{\ss}ler,
Y.~Matiyasevich and D.~Schleicher, \emph{Diophantine equations over
\(\mathbb Z\): universal bounds and parallel formalization},
arXiv:2506.20909v3, Definitions 2.1, 4.1 and 5.3;
\url{https://arxiv.org/abs/2506.20909v3}.

\begin{proposition}
For every integer-coefficient input polynomial \(P\) and every
natural-number instance parameter \(a\), the constructed source
\(Q\) has a solution with all nine inputs integral and \(n=-7\);
it also has one with \(n=-9\).  Consequently the other eight
integrality conditions do not make \(n\ge0\) or \(R>0\) automatic.
\end{proposition}
\begin{proof}
Choose
\[
                f=g=l=x=0,\qquad w=2,\qquad y=0.
\]
Then \(b=1\), and \(B_0=\beta\) is divisible by four.
The constant digit of the mask is \(B_0-1\), whereas all higher
digits multiply positive powers of \(B_0\).  Thus the mask is odd,
and \(T_0\) is odd.  Also \(N\) is even.  It follows that
\[
 R_0=(S_0+T_0+1)N+T_0+1
\]
is even, and hence \(X=(N-1)R_0\) is even.  Therefore the
additional choices
\[
                         h=X,\qquad k=X/2
\]
are integers.

The bridge definitions now give, by direct substitution,
\[
 U=V=A=0,\quad B_1=2X+1,\quad C=J=1,\quad D=E=0,
 \quad F=G=I=1.
\]
In particular,
\[
        A_1=1,\qquad A_2=A_3=0,\qquad
        S=-5,\quad T=R=0,\quad Z_0=2.
\]
Every sign factor defining \(Q\) is therefore
\(25(n+8+\varepsilon_1)\).  The other two signs have no
effect, so exactly
\[
                     Q(n)=25^8(n+9)^4(n+7)^4.
\]
Both displayed negative inputs are roots, independently of the
input polynomial and the instance parameter.
\end{proof}

This explicit obstruction does not contradict the published source
theorem, which requires \(n\ge0\).  It explains why the finite-wrapper
transfer retains both the ninth integer predicate and a genuine
nonnegative representation.  The construction also shows that merely
requiring all nine inputs to be integers cannot replace that wrapper.
\end{document}
